% =====================================================================
%  Appendix A -- Proofs for Sections 3-4 (owner: THEORY).
% =====================================================================
\providecommand{\osc}{\operatorname{osc}}
\providecommand{\dTV}{d_{\mathrm{TV}}}

\section{Proofs and further results for Sections~\ref{sec:estimands}--\ref{sec:markov}}\label{app:proofs-estimands}

Throughout, vectors are columns, $\bm1$ is the vector of ones, $\bm e_x$ is the $x$th unit vector, and for a matrix $B$ we write $B_{x\cdot}$ for its $x$th row, as a row vector. Two elementary lemmas are used repeatedly.

\begin{lemma}[Zero-sum vectors]\label{lem:app-zerosum}
Let $\bm v\in\mathbb R^{\mathcal X}$ with $\bm v^{\top}\bm1=0$.
\begin{enumerate}[(a)]
  \item For every $\bm h\in\mathbb R^m$, $|\bm v^{\top}\bm h|\le\frac12\|\bm v\|_1\osc(\bm h)$.
  \item For every real matrix $B$ with rows indexed by $\mathcal X$, $\|\bm v^{\top}B\|_1\le\operatorname{erg}(B)\,\|\bm v\|_1$, where $\operatorname{erg}(B)=\frac12\max_{x,y}\|B_{x\cdot}-B_{y\cdot}\|_1$.
  \item If every row of $B$ has the same sum, then $\osc(B\bm h)\le\operatorname{erg}(B)\osc(\bm h)$.
\end{enumerate}
\end{lemma}

\begin{proof}
Write $\bm v=\bm v^+-\bm v^-$ with $\bm v^\pm\ge\bm0$ the positive and negative parts. Since $\bm v^{\top}\bm1=0$, both parts have total mass $s=\frac12\|\bm v\|_1$.
(a) $\bm v^{\top}\bm h=(\bm v^+)^{\top}\bm h-(\bm v^-)^{\top}\bm h\le s\max\bm h-s\min\bm h$, and the same bound holds for $-\bm v$.
(b) If $s=0$ there is nothing to prove. Otherwise $\bm p=\bm v^+/s$ and $\bm q=\bm v^-/s$ are probability vectors and $\bm v^{\top}B=s\sum_{x,y}p_xq_y(B_{x\cdot}-B_{y\cdot})$. By the triangle inequality, $\|\bm v^{\top}B\|_1\le s\max_{x,y}\|B_{x\cdot}-B_{y\cdot}\|_1=2s\,\operatorname{erg}(B)$.
(c) For rows $x,y$, the vector $(B_{x\cdot}-B_{y\cdot})^{\top}$ sums to zero, so (a) gives $(B\bm h)_x-(B\bm h)_y\le\frac12\|B_{x\cdot}-B_{y\cdot}\|_1\osc(\bm h)\le\operatorname{erg}(B)\osc(\bm h)$.
\end{proof}

\begin{lemma}[Integrals against a difference of distributions]\label{lem:app-signed}
Let $F$ and $G$ be probability measures on a measurable space $\mathcal D$, let $S$ be a measurable set with $F(S)=G(S)=1$, and let $f$ be bounded and measurable. Then $|\int f\,d(F-G)|\le\osc_S(f)\,\dTV(F,G)$, where $\dTV(F,G)=\sup_B|F(B)-G(B)|$. If $\mathcal D$ is a metric space, $F$ and $G$ have finite first moments and $f$ is $L$-Lipschitz, then $|\int f\,d(F-G)|\le L\,\dW(F,G)$.
\end{lemma}

\begin{proof}
Write the Jordan decomposition of the signed measure $F-G$ as $(F-G)^+-(F-G)^-$. Since $(F-G)(\mathcal D)=0$, $(F-G)^+(\mathcal D)=(F-G)^-(\mathcal D)=\sup_B(F-G)(B)=\dTV(F,G)$. Both parts vanish outside $S$, because $|F-G|(\mathcal D\setminus S)\le F(\mathcal D\setminus S)+G(\mathcal D\setminus S)=0$. Hence $\int f\,d(F-G)=\int_Sf\,d(F-G)^+-\int_Sf\,d(F-G)^-\le\dTV(F,G)(\sup_Sf-\inf_Sf)$, and the same holds for $-f$. The second bound is the Kantorovich--Rubinstein duality, $\dW(F,G)=\sup\{\int h\,d(F-G):h\text{ $1$-Lipschitz}\}$ \citep[Ch.~6]{villani2009}, applied to $h=f/L$.
\end{proof}

\subsection{Proof of \texorpdfstring{\cref{thm:decomp}}{the decomposition theorem}}\label{app:pf-decomp}

\emph{(a)} By the definitions \eqref{eq:est-components}, $\tau_{\mathrm{sup}}+\tau_{\mathrm{dem}}+\iota=\bar\mu(\mathrm{alg};\mathrm{alg})-\bar\mu(\mathrm{ref};\mathrm{ref})=\tau$, which is the first identity. The two paths are the telescoping sums
\[
\begin{aligned}
  \tau&=\underbrace{\{\bar\mu(\mathrm{alg};\mathrm{ref})-\bar\mu(\mathrm{ref};\mathrm{ref})\}}_{\tau_{\mathrm{sup}}}
  +\underbrace{\{\bar\mu(\mathrm{alg};\mathrm{alg})-\bar\mu(\mathrm{alg};\mathrm{ref})\}}_{\tau_{\mathrm{dem}}+\iota}\\
  &=\underbrace{\{\bar\mu(\mathrm{alg};\mathrm{alg})-\bar\mu(\mathrm{ref};\mathrm{alg})\}}_{\tau_{\mathrm{sup}}+\iota}
  +\underbrace{\{\bar\mu(\mathrm{ref};\mathrm{alg})-\bar\mu(\mathrm{ref};\mathrm{ref})\}}_{\tau_{\mathrm{dem}}},
\end{aligned}
\]
where the labels under the braces follow from the definition of $\iota$.

\emph{(b)} The game has $v(\emptyset)=0$, $v(\{\mathrm{sup}\})=\tau_{\mathrm{sup}}$, $v(\{\mathrm{dem}\})=\tau_{\mathrm{dem}}$ and $v(\{\mathrm{sup},\mathrm{dem}\})=\tau$. With two players, the Shapley value averages each player's marginal contribution over the two orders \citep{shapley1953}:
$\phi_{\mathrm{sup}}=\frac12\{v(\{\mathrm{sup}\})-v(\emptyset)\}+\frac12\{v(\{\mathrm{sup},\mathrm{dem}\})-v(\{\mathrm{dem}\})\}=\frac12\tau_{\mathrm{sup}}+\frac12(\tau_{\mathrm{sup}}+\iota)$. This is the average of the two path attributions of (a), and it equals $\tau_{\mathrm{sup}}+\iota/2$. The expression for $\phi_{\mathrm{dem}}$ follows by symmetry, and $\phi_{\mathrm{sup}}+\phi_{\mathrm{dem}}=\tau_{\mathrm{sup}}+\tau_{\mathrm{dem}}+\iota=\tau$.
For uniqueness, write $v=v_1+v_2+v_{12}$ with $v_1(J)=\tau_{\mathrm{sup}}\ind{\mathrm{sup}\in J}$, $v_2(J)=\tau_{\mathrm{dem}}\ind{\mathrm{dem}\in J}$ and $v_{12}(J)=\iota\,\ind{J=\{\mathrm{sup},\mathrm{dem}\}}$. Let $\psi$ be an allocation that is efficient, symmetric, additive and gives nothing to null players. In $v_1$ demand is a null player, so $\psi_{\mathrm{dem}}(v_1)=0$ and, by efficiency, $\psi_{\mathrm{sup}}(v_1)=\tau_{\mathrm{sup}}$. Likewise $\psi(v_2)=(0,\tau_{\mathrm{dem}})$. In $v_{12}$ the players are symmetric, so each receives $\iota/2$. Additivity gives $\psi(v)=(\tau_{\mathrm{sup}}+\iota/2,\tau_{\mathrm{dem}}+\iota/2)=\phi$.
For the ratio scale, apply the same computation to $v^{\log}(J)=\log\bar\mu(a_J;b_J)-\log\bar\mu(\mathrm{ref};\mathrm{ref})$. Then
\[
  \phi^{\log}_{\mathrm{sup}}=\tfrac12\log\frac{\bar\mu(\mathrm{alg};\mathrm{ref})}{\bar\mu(\mathrm{ref};\mathrm{ref})}+\tfrac12\log\frac{\bar\mu(\mathrm{alg};\mathrm{alg})}{\bar\mu(\mathrm{ref};\mathrm{alg})}
  =\tfrac12\log\{\mathrm{AR}(F_{\mathrm{ref}})\,\mathrm{AR}(F_{\mathrm{alg}})\},
\]
and similarly $\phi^{\log}_{\mathrm{dem}}=\frac12\log\{\mathrm{DR}(\mathrm{ref})\,\mathrm{DR}(\mathrm{alg})\}$. Efficiency gives $\phi^{\log}_{\mathrm{sup}}+\phi^{\log}_{\mathrm{dem}}=\log\mathrm{TR}$, and exponentiating gives $\mathrm{TR}=\Phi_{\mathrm{sup}}\Phi_{\mathrm{dem}}$. The log interaction is
\[
  v^{\log}(\{\mathrm{sup},\mathrm{dem}\})-v^{\log}(\{\mathrm{sup}\})-v^{\log}(\{\mathrm{dem}\})
  =\log\frac{\bar\mu(\mathrm{alg};\mathrm{alg})\,\bar\mu(\mathrm{ref};\mathrm{ref})}{\bar\mu(\mathrm{alg};\mathrm{ref})\,\bar\mu(\mathrm{ref};\mathrm{alg})},
\]
whose exponential can be grouped either as $\mathrm{AR}(F_{\mathrm{alg}})/\mathrm{AR}(F_{\mathrm{ref}})$ or as $\mathrm{DR}(\mathrm{alg})/\mathrm{DR}(\mathrm{ref})$.

\emph{(c)} By \cref{asm:demand} and iterated expectations, $\bar\mu(a;b)=\E[\E\{E_U(a;b)\mid D_U(b)\}]=\E\{\mu_a(D_U(b))\}=\int\mu_a\,dF_b$; the integrals exist because $0\le\mu_a\le1$. This is \eqref{eq:est-bilinear}. Substituting it into \eqref{eq:est-components} gives $\tau_{\mathrm{sup}}=\int\Delta\mu\,dF_{\mathrm{ref}}=\mathrm{AA}_{\mathrm{ref}}(F_{\mathrm{ref}})$, $\tau_{\mathrm{dem}}=\int\mu_{\mathrm{ref}}\,d(F_{\mathrm{alg}}-F_{\mathrm{ref}})$ and
\[
  \iota=\int\mu_{\mathrm{alg}}\,dF_{\mathrm{alg}}-\int\mu_{\mathrm{alg}}\,dF_{\mathrm{ref}}-\int\mu_{\mathrm{ref}}\,dF_{\mathrm{alg}}+\int\mu_{\mathrm{ref}}\,dF_{\mathrm{ref}}=\int\Delta\mu\,d(F_{\mathrm{alg}}-F_{\mathrm{ref}}).
\]
Hence $\tau_{\mathrm{sup}}+\iota=\int\Delta\mu\,dF_{\mathrm{alg}}=\mathrm{AA}_{\mathrm{ref}}(F_{\mathrm{alg}})$. Finally, $\phi_{\mathrm{sup}}=\tau_{\mathrm{sup}}+\frac12\iota=\frac12\int\Delta\mu\,dF_{\mathrm{ref}}+\frac12\int\Delta\mu\,dF_{\mathrm{alg}}=\int\Delta\mu\,dF_{\mathrm{mid}}$, and
$\phi_{\mathrm{dem}}=\tau_{\mathrm{dem}}+\frac12\iota=\int\{\mu_{\mathrm{ref}}+\frac12(\mu_{\mathrm{alg}}-\mu_{\mathrm{ref}})\}\,d(F_{\mathrm{alg}}-F_{\mathrm{ref}})=\int\mu_{\mathrm{mid}}\,d(F_{\mathrm{alg}}-F_{\mathrm{ref}})$.

\emph{(d)} By (c), $\iota=\mathrm{AA}_{\mathrm{ref}}(F_{\mathrm{alg}})-\mathrm{AA}_{\mathrm{ref}}(F_{\mathrm{ref}})$, which vanishes exactly when the two are equal. By (b) and \eqref{eq:est-bilinear}, $\iota^{\times}=\mathrm{AR}_{\mathrm{ref}}(F_{\mathrm{alg}})/\mathrm{AR}_{\mathrm{ref}}(F_{\mathrm{ref}})$, which equals one exactly when the two ratios are equal.
In case (i), every integral against $F_{\mathrm{alg}}-F_{\mathrm{ref}}$ vanishes, so $\iota=\tau_{\mathrm{dem}}=0$, the two path attributions in (a) both equal $\tau_{\mathrm{sup}}=\tau$, and $\bar\mu(\mathrm{alg};\mathrm{alg})=\mu_{\mathrm{alg}}(F)=\mu_{\mathrm{ref}}(F)+\mathrm{AA}_{\mathrm{ref}}(F)$.
In case (ii), $\Delta\mu=c$ on $S$, and $F_{\mathrm{alg}}$ and $F_{\mathrm{ref}}$ are both concentrated on $S$, so $\iota=c\{F_{\mathrm{alg}}(S)-F_{\mathrm{ref}}(S)\}=0$.
The bounds follow from \cref{lem:app-signed}, applied to $f=\Delta\mu$ for $\iota$ and to $f=\mu_{\mathrm{ref}}$ for $\tau_{\mathrm{dem}}$. \qed

\subsection{Proof of \texorpdfstring{\cref{prop:audit-scope}}{the audit-scope proposition}}\label{app:pf-audit}

\emph{(a)} $\mathrm{AR}_\rho(\Lambda)$ and $\mathrm{AA}_\rho(\Lambda)$ are functionals of $\Lambda$ and of $\mu_{\mathrm{alg}}$ and $\mu_\rho$ restricted to a set of $\Lambda$-probability one, here $\mathcal D_A$. The three special cases are \cref{thm:decomp}(c).

\emph{(b)} The audit observes only sessions of puppets whose demand is fixed at personas in $\mathcal D_A$. Under the conditions of \cref{sec:identification}, the joint law of these sessions is determined by the policies and the personas. It does not involve the demand profiles $D_U(b)$ of real users, so it is the same for every specification of the formation laws $(F_{\mathrm{alg}},F_{\mathrm{ref}})$. Grant oracle knowledge: fix $\mu_{\mathrm{alg}}$ and $\mu_{\mathrm{ref}}$ on all of $\mathcal D$, and fix $F_{\mathrm{alg}}$. For an arbitrary probability measure $G$ on $\mathcal D$, a population model with $D_U(\mathrm{ref})\sim G$ and $\E\{E_U(a;\mathrm{ref})\mid D_U(\mathrm{ref})=d\}=\mu_a(d)$ satisfies \cref{asm:demand}. It leaves the law of the audit data unchanged and has $F_{\mathrm{ref}}=G$. By \eqref{eq:est-bilinear}, $\tau=\bar\mu(\mathrm{alg};\mathrm{alg})-\mu_{\mathrm{ref}}(G)$, where $\bar\mu(\mathrm{alg};\mathrm{alg})=\mu_{\mathrm{alg}}(F_{\mathrm{alg}})$ is fixed. The identified set for $\tau$ is therefore $\{\bar\mu(\mathrm{alg};\mathrm{alg})-\mu_{\mathrm{ref}}(G)\}$, as $G$ ranges over all probability measures on $\mathcal D$.
It remains to describe $\{\mu_{\mathrm{ref}}(G)\}$. Let $s=\sup_{\mathcal D}\mu_{\mathrm{ref}}$ and $i=\inf_{\mathcal D}\mu_{\mathrm{ref}}$. Every $\mu_{\mathrm{ref}}(G)$ lies in $[i,s]$. Point masses give every value $\mu_{\mathrm{ref}}(d)$, and two-point mixtures give every convex combination of two such values, so the set contains the open interval $(i,s)$. The endpoint $s$ belongs to the set if and only if it is attained: if $\mu_{\mathrm{ref}}<s$ everywhere, then $\mu_{\mathrm{ref}}(G)<s$ for every $G$. The same holds for $i$. The identified set for $\tau$ is thus an interval with endpoints $\bar\mu(\mathrm{alg};\mathrm{alg})-s$ and $\bar\mu(\mathrm{alg};\mathrm{alg})-i$, of length $s-i=\osc_{\mathcal D}(\mu_{\mathrm{ref}})$. Since $\tau_{\mathrm{dem}}=\mu_{\mathrm{ref}}(F_{\mathrm{alg}})-\mu_{\mathrm{ref}}(G)$, the same argument applies to $\tau_{\mathrm{dem}}$. In particular, $\tau$ and $\tau_{\mathrm{dem}}$ are identified if and only if this interval is a point, that is, if and only if $\mu_{\mathrm{ref}}$ is constant on $\mathcal D$. Likewise, as $G$ ranges over the point masses, $\iota=\mathrm{AA}_{\mathrm{ref}}(F_{\mathrm{alg}})-\int\Delta\mu\,dG$ and $\phi_{\mathrm{dem}}=\int\mu_{\mathrm{mid}}\,d(F_{\mathrm{alg}}-G)$ take every value of $\mathrm{AA}_{\mathrm{ref}}(F_{\mathrm{alg}})-\Delta\mu(d)$ and $\mu_{\mathrm{mid}}(F_{\mathrm{alg}})-\mu_{\mathrm{mid}}(d)$, so they are identified if and only if $\Delta\mu$, respectively $\mu_{\mathrm{mid}}$, is constant on $\mathcal D$. Without oracle knowledge, the values of $\mu_{\mathrm{alg}}$ and $\mu_{\mathrm{ref}}$ off $\mathcal D_A$ and the law $F_{\mathrm{alg}}$ are further unknowns, so the identified sets can only be larger.

\emph{(c)} By \eqref{eq:est-components}, $\tau-\mathrm{AA}_{\mathrm{ref}}(F_{\mathrm{alg}})=\tau-(\tau_{\mathrm{sup}}+\iota)=\tau_{\mathrm{dem}}$ and $\phi_{\mathrm{sup}}-(\tau_{\mathrm{sup}}+\iota)=-\iota/2$. The bounds follow from \cref{thm:decomp}(d). \qed

\subsection{Proof of \texorpdfstring{\cref{prop:markov}}{the exposure-functionals proposition}}\label{app:pf-markov}

We use three standard facts about a unichain $P$ with closed class $C$ \citep[Chs.~3--5]{kemeny1960}; see also \citet{seneta2006}. (U1) $P$ has a unique stationary distribution $\varpi$, and $\varpi_x>0$ if and only if $x\in C$. (U2) $T^{-1}\sum_{t=0}^{T-1}P^t\to\bm1\varpi^{\top}$. (U3) Every solution of $P\bm v=\bm v$ is constant, because $\bm v=T^{-1}\sum_{t<T}P^t\bm v\to\bm1\varpi^{\top}\bm v$.

\emph{(a) Algebra.} If $(I-P+\bm1\varpi^{\top})\bm v=\bm0$, multiplying on the left by $\varpi^{\top}$ gives $\varpi^{\top}\bm v=0$. Then $(I-P)\bm v=\bm0$, so $\bm v$ is constant by (U3), and $\bm v=\bm0$. Thus $Z$ exists. Since $(I-P+\bm1\varpi^{\top})\bm1=\bm1$ and $\varpi^{\top}(I-P+\bm1\varpi^{\top})=\varpi^{\top}$, we have $Z\bm1=\bm1$ and $\varpi^{\top}Z=\varpi^{\top}$. Hence $(I-P)Z=(Z^{-1}-\bm1\varpi^{\top})Z=I-\bm1\varpi^{\top}$, and likewise $Z(I-P)=I-\bm1\varpi^{\top}$. For $H=Z-\bm1\varpi^{\top}$ this gives $H\bm1=\bm0$, $\varpi^{\top}H=\bm0^{\top}$ and $(I-P)H=H(I-P)=I-\bm1\varpi^{\top}$. Therefore $H(I-P)H=H-\bm1\varpi^{\top}H=H$ and $(I-P)H(I-P)=(I-\bm1\varpi^{\top})(I-P)=I-P$. These are the defining equations of the group inverse \citep{meyer1975}. For the formula for $\varpi$: $\varpi^{\top}(I-P+\bm1\bm1^{\top})=\bm1^{\top}$, and $I-P+\bm1\bm1^{\top}$ is invertible by the argument used for $Z$. If $(I-P+\bm1\bm1^{\top})\bm v=\bm0$, then multiplying by $\varpi^{\top}$ gives $\bm1^{\top}\bm v=0$, so $(I-P)\bm v=\bm0$, $\bm v=c\bm1$ and $c|\mathcal X|=0$.

\emph{(a) Limit theorems.} Let $\tilde{\bm f}=\bm1_{\mathcal A}-\varpi(\mathcal A)\bm1$ and $\bm g=\bm g_{\mathcal A}=H\bm1_{\mathcal A}=H\tilde{\bm f}$. Then $(I-P)\bm g=(I-\bm1\varpi^{\top})\bm1_{\mathcal A}=\tilde{\bm f}$, so $\tilde f(X_t)=g(X_t)-(P\bm g)(X_t)$ and
\begin{equation}\label{eq:app-poisson}
  \sum_{t=1}^{T}\tilde f(X_t)=\sum_{t=1}^{T-1}d_{t+1}+g(X_1)-(P\bm g)(X_T),\qquad d_{t+1}=g(X_{t+1})-(P\bm g)(X_t).
\end{equation}
The $d_{t+1}$ are bounded martingale differences for the natural filtration, with conditional variances $\sigma^2_g(X_t)$, where $\sigma^2_g(x)=\Var_{y\sim P(x,\cdot)}\{g(y)\}$. By the strong law for martingales with bounded increments \citep[Ch.~2]{hall1980}, $T^{-1}\sum_td_{t+1}\to0$ almost surely, and the boundary terms in \eqref{eq:app-poisson} are bounded. Hence $E-\varpi(\mathcal A)=T^{-1}\sum_t\tilde f(X_t)\to0$ almost surely. The same argument applied to $\sigma^2_g$ instead of $\bm1_{\mathcal A}$ gives $T^{-1}\sum_t\sigma^2_g(X_t)\to\sum_x\varpi_x\sigma^2_g(x)$ almost surely. The martingale central limit theorem \citep[Ch.~3]{hall1980} then yields $T^{-1/2}\sum_td_{t+1}\Rightarrow\mathrm N(0,\sum_x\varpi_x\sigma^2_g(x))$, and the boundary terms are negligible after division by $\sqrt T$. To identify the variance, use $\varpi^{\top}P=\varpi^{\top}$ and $P\bm g=\bm g-\tilde{\bm f}$:
\[
\begin{aligned}
  \sum_x\varpi_x\sigma^2_g(x)&=\varpi^{\top}P(\bm g\circ\bm g)-\varpi^{\top}\{(P\bm g)\circ(P\bm g)\}
  =\varpi^{\top}(\bm g\circ\bm g)-\varpi^{\top}\{(\bm g-\tilde{\bm f})\circ(\bm g-\tilde{\bm f})\}\\
  &=2\varpi^{\top}(\tilde{\bm f}\circ\bm g)-\varpi^{\top}(\tilde{\bm f}\circ\tilde{\bm f}),
\end{aligned}
\]
where $\circ$ is the entrywise product. Since $\bm g=H\tilde{\bm f}$, this is the first expression for $\sigma^2_{\mathcal A}$.

\emph{(b)} Multiplying the telescoping identity $\sum_{t=0}^{T-1}P^t(I-P)=I-P^T$ on the right by $H$, and using $(I-P)H=I-\bm1\varpi^{\top}$ and $P^t\bm1=\bm1$, gives
\begin{equation}\label{eq:app-sumP}
  \sum_{t=0}^{T-1}P^t=T\,\bm1\varpi^{\top}+(I-P^T)H .
\end{equation}
Substituting into \eqref{eq:mk-eT} gives $e_T=\varpi(\mathcal A)+T^{-1}\nu^{\top}(I-P^T)\bm g_{\mathcal A}$. Since $(I-P^T)\bm1=\bm0$, $(I-P^T)H=(I-P^T)Z$, which gives the second form. Write $\nu^{\top}(I-P^T)\bm g_{\mathcal A}=b(\nu)+r_T$ with $b(\nu)=\nu^{\top}\bm g_{\mathcal A}$ and $r_T=-\nu^{\top}P^T\bm g_{\mathcal A}=-(\nu^{\top}P^T-\varpi^{\top})\bm g_{\mathcal A}$, using $\varpi^{\top}\bm g_{\mathcal A}=0$. Also $b(\nu)=(\nu-\varpi)^{\top}\bm g_{\mathcal A}=(\nu-\varpi)^{\top}Z\bm1_{\mathcal A}$, because $(\nu-\varpi)^{\top}\bm1=0$. The vector $\nu^{\top}P^T-\varpi^{\top}$ sums to zero, so \cref{lem:app-zerosum}(a) bounds $|r_T|$ as stated. If $P$ is aperiodic, then $\nu^{\top}P^T\to\varpi^{\top}$.

\emph{(c)} The event that a session has no antisemitic item is $\{X_1,\dots,X_T\in\bar{\mathcal A}\}$, whose probability is $\sum\nu_{x_1}P_{x_1x_2}\cdots P_{x_{T-1}x_T}$ over $x_1,\dots,x_T\in\bar{\mathcal A}$, that is, $\nu_{\bar{\mathcal A}}^{\top}P_{\bar{\mathcal A}\bar{\mathcal A}}^{T-1}\bm1$. If $\mathcal A$ is accessible from every state, then making the states of $\mathcal A$ absorbing gives an absorbing chain, so the spectral radius of $P_{\bar{\mathcal A}\bar{\mathcal A}}$ is less than one \citep[Ch.~3]{kemeny1960}, and $I-P_{\bar{\mathcal A}\bar{\mathcal A}}$ is invertible. For $x\in\bar{\mathcal A}$, $\Pr_x(T_{\mathcal A}>n)=\Pr_x(X_1,\dots,X_{1+n}\in\bar{\mathcal A})=(P_{\bar{\mathcal A}\bar{\mathcal A}}^{n}\bm1)_x$, so $M_{x\mathcal A}=\sum_{n\ge0}(P_{\bar{\mathcal A}\bar{\mathcal A}}^n\bm1)_x=\{(I-P_{\bar{\mathcal A}\bar{\mathcal A}})^{-1}\bm1\}_x$.
Now let $P$ be irreducible and fix $y$, so that $\varpi_y>0$. Let $u_x=(H_{yy}-H_{xy})/\varpi_y$. Then $u_y=0$. The $y$th column of $(I-P)H=I-\bm1\varpi^{\top}$ reads $(I-P)H\bm e_y=\bm e_y-\varpi_y\bm1$, so $(I-P)\bm u=-(I-P)H\bm e_y/\varpi_y=\bm1-\bm e_y/\varpi_y$. For $x\neq y$ this says $u_x-\sum_{z\neq y}P_{xz}u_z=1$, using $u_y=0$. With $\mathcal X_{-y}=\mathcal X\setminus\{y\}$, this is the system $(I-P_{\mathcal X_{-y}\mathcal X_{-y}})\bm u_{\mathcal X_{-y}}=\bm1$, whose unique solution is $(M_{xy})_{x\neq y}$. Hence $M_{xy}=u_x$ for all $x$. Summing, $\sum_y\varpi_yM_{xy}=\sum_y(H_{yy}-H_{xy})=\operatorname{tr}H-(H\bm1)_x=\operatorname{tr}H$.
For the spectrum, let $\mathcal R=\{\bm v:\varpi^{\top}\bm v=0\}$. The space $\mathbb R^{\mathcal X}$ is the direct sum of $\operatorname{span}(\bm1)$ and $\mathcal R$, and both are invariant under $P$ and under $H$, because $P\bm1=\bm1$, $H\bm1=\bm0$, $\varpi^{\top}P=\varpi^{\top}$ and $\varpi^{\top}H=\bm0^{\top}$. On $\mathcal R$, $(I-P)H$ and $H(I-P)$ act as the identity, because $I-\bm1\varpi^{\top}$ does. So $H$ restricted to $\mathcal R$ is the inverse of $I-P$ restricted to $\mathcal R$, and $H$ is zero on $\operatorname{span}(\bm1)$. The eigenvalues of $P$ restricted to $\mathcal R$ are the eigenvalues of $P$ with one copy of the eigenvalue $1$ removed, namely $\lambda_2,\dots,\lambda_{|\mathcal X|}$, none of which equals one because $1$ is a simple eigenvalue of an irreducible $P$. Hence $\operatorname{tr}H=\sum_{k=2}^{|\mathcal X|}(1-\lambda_k)^{-1}$.

\emph{(d)} From every transient state some closed class is accessible, so, as in (c), the spectral radius of $P_{\mathcal T\mathcal T}$ is less than one and $Z_{\mathcal T}=\sum_{n\ge0}P_{\mathcal T\mathcal T}^n$ exists. The probability that a session started at $x\in\mathcal T$ stays in $\mathcal T$ for items $1,\dots,n+1$ and enters $C_k$ at item $n+2$ is $(P_{\mathcal T\mathcal T}^nP_{\mathcal TC_k}\bm1)_x$, so $B_{xk}=\sum_{n\ge0}(P_{\mathcal T\mathcal T}^nP_{\mathcal TC_k}\bm1)_x=(Z_{\mathcal T}P_{\mathcal TC_k}\bm1)_x$, and the expected number of items before entering a closed class is $\sum_n\Pr_x(X_1,\dots,X_{1+n}\in\mathcal T)=(Z_{\mathcal T}\bm1)_x$. For the long-run exposure, let $T_{\mathcal C}$ be the first time the chain is in a closed class. It is almost surely finite. On the event that the chain enters $C_k$, the strong Markov property and part (a), applied to the irreducible chain $P$ restricted to $C_k$, give $T^{-1}\sum_{t=T_{\mathcal C}}^{T}\ind{X_t\in\mathcal A}\to\varpi^{(k)}(\mathcal A)$ almost surely, while $T^{-1}\sum_{t<T_{\mathcal C}}\ind{X_t\in\mathcal A}\le T_{\mathcal C}/T\to0$. Bounded convergence gives $\lim_Te_T(P,\nu)=\sum_k\Pr_\nu(\text{enter }C_k)\,\varpi^{(k)}(\mathcal A)$, and $\Pr_\nu(\text{enter }C_k)=\nu(C_k)+\nu_{\mathcal T}^{\top}Z_{\mathcal T}P_{\mathcal TC_k}\bm1$. \qed

\subsection{Proof of \texorpdfstring{\cref{thm:perturb}}{the sensitivity theorem}}\label{app:pf-perturb}

Write $P=P_{\mathrm{ref}}$, $\tilde P=P_{\mathrm{alg}}$, $\tilde\varpi=\varpi_{\mathrm{alg}}$ and $\Delta P=\tilde P-P$, so that $\Delta P\bm1=\bm0$, and $\bm g=\bm g_{\mathcal A}$.

\emph{(a)} Since $\tilde\varpi^{\top}\tilde P=\tilde\varpi^{\top}$, $\tilde\varpi^{\top}(I-P)=\tilde\varpi^{\top}(I-\tilde P)+\tilde\varpi^{\top}\Delta P=\tilde\varpi^{\top}\Delta P$. Multiplying on the right by $H$ and using $(I-P)H=I-\bm1\varpi^{\top}$ and $\tilde\varpi^{\top}\bm1=1$ gives $\tilde\varpi^{\top}-\varpi^{\top}=\tilde\varpi^{\top}\Delta P H$, so $\tilde\varpi^{\top}(I-\Delta P H)=\varpi^{\top}$. To see that $I-\Delta P H$ is invertible, suppose $\bm v^{\top}(I-\Delta P H)=\bm0^{\top}$. Then $\bm v^{\top}\bm1=\bm v^{\top}\Delta P H\bm1=0$, and $\bm v^{\top}(I-P)=\bm v^{\top}\Delta P H(I-P)=\bm v^{\top}\Delta P(I-\bm1\varpi^{\top})=\bm v^{\top}\Delta P$. Hence $\bm v^{\top}(I-\tilde P)=\bm0^{\top}$. Because $\tilde P$ is unichain, $\bm v=c\tilde\varpi$, and $c=\bm v^{\top}\bm1=0$. Thus $\tilde\varpi^{\top}=\varpi^{\top}(I-\Delta P H)^{-1}$.
Along the path, apply this with $\Delta P=\varepsilon\Delta$: $\varpi(\varepsilon)^{\top}=\varpi^{\top}(I-\varepsilon\Delta H)^{-1}$, which is rational in $\varepsilon$ by Cramer's rule and continuous wherever $P(\varepsilon)$ is unichain. The identity only requires the base matrix to be unichain, so it can be applied at any point $\varepsilon$ of the path with increment $h\Delta$: $\varpi(\varepsilon+h)^{\top}-\varpi(\varepsilon)^{\top}=h\,\varpi(\varepsilon+h)^{\top}\Delta H(\varepsilon)$. Dividing by $h$ and letting $h\to0$ gives $d\varpi(\varepsilon)^{\top}/d\varepsilon=\varpi(\varepsilon)^{\top}\Delta H(\varepsilon)$. Since $Z(\varepsilon)=H(\varepsilon)+\bm1\varpi(\varepsilon)^{\top}$ and $\Delta\bm1=\bm0$, $\Delta Z(\varepsilon)=\Delta H(\varepsilon)$. Multiplying by $\bm1_{\mathcal A}$ at $\varepsilon=0$ gives the derivative of $\varpi_{\mathcal A}$.

\emph{(b)} By (a), $\tilde\varpi(\mathcal A)-\varpi(\mathcal A)=\bm v^{\top}\bm g$ and $\tilde\varpi^{\top}-\varpi^{\top}=\bm v^{\top}H$, with $\bm v^{\top}=\tilde\varpi^{\top}\Delta P$. Then $\bm v^{\top}\bm1=0$ and $\|\bm v\|_1\le\sum_x\tilde\varpi_x\|\Delta P_{x\cdot}\|_1\le\|\Delta P\|_\infty$. \Cref{lem:app-zerosum}(a) gives $|\bm v^{\top}\bm g|\le\frac12\|\Delta P\|_\infty\osc(\bm g)=\kappa_{\mathcal A}\|\Delta P\|_\infty$, and \cref{lem:app-zerosum}(b) gives $\|\bm v^{\top}H\|_1\le\operatorname{erg}(H)\|\Delta P\|_\infty$.
By \cref{prop:markov}(c), $H_{xy}=H_{yy}-\varpi_yM_{xy}$, so $g_x=\sum_{y\in\mathcal A}H_{yy}-\sum_{y\in\mathcal A}\varpi_yM_{xy}$ and $\osc(\bm g)=\osc_x\sum_{y\in\mathcal A}\varpi_yM_{xy}$. Each sum lies in $[0,\sum_y\varpi_yM_{xy}]=[0,K(P)]$, so $\osc(\bm g)\le K(P)$. For any $x,z$, the vector $(H_{x\cdot}-H_{z\cdot})^{\top}$ sums to zero, so by \cref{lem:app-zerosum}(a), $g_x-g_z=(H_{x\cdot}-H_{z\cdot})\bm1_{\mathcal A}\le\frac12\|H_{x\cdot}-H_{z\cdot}\|_1\le\operatorname{erg}(H)$. This proves $\kappa_{\mathcal A}\le\frac12\min\{\operatorname{erg}(H),K(P)\}$. For a single state $y$, $\varpi_yM_{xy}$ is minimised at $x=y$, where it is zero, so $\kappa_y=\frac12\varpi_y\max_xM_{xy}$.
For sharpness, let $\Delta\bm1=\bm0$ and $\|\Delta\|_\infty\le1$. By \cref{lem:app-zerosum}(a) applied row by row, $|\varpi^{\top}\Delta\bm g|\le\sum_x\varpi_x\frac12\|\Delta_{x\cdot}\|_1\osc(\bm g)\le\kappa_{\mathcal A}$. Let $x^+$ maximise and $x^-$ minimise $\bm g$, and let every row of $\Delta$ be $\frac12(\bm e_{x^+}-\bm e_{x^-})^{\top}$. Then $\|\Delta\|_\infty=1$, $\Delta_{x\cdot}\bm g=\frac12\osc(\bm g)$ for every $x$, and $\varpi^{\top}\Delta\bm g=\kappa_{\mathcal A}$. The matrix $P+\varepsilon\Delta$ is stochastic for small $\varepsilon>0$ if and only if $P_{xx^-}>0$ for every $x$. In general, a direction keeps $P+\varepsilon\Delta$ stochastic for small $\varepsilon>0$ if and only if the negative entries of each row $\Delta_{x\cdot}$ lie in the support $S_x=\{y:P_{xy}>0\}$. Then, by the argument of \cref{lem:app-zerosum}(a), $\Delta_{x\cdot}\bm g\le\frac12\|\Delta_{x\cdot}\|_1\{\max_yg_y-\min_{y\in S_x}g_y\}$, with equality for the row $\frac12(\bm e_{x^+}-\bm e_{y_x})^{\top}$, where $y_x$ minimises $\bm g$ over $S_x$. Weighting by $\varpi_x$ and summing gives the largest first-order increase stated in \cref{rem:perturb-sharp}.

\emph{(c)} The matrix $P(\varepsilon)$ is stochastic and, for $\varepsilon>0$, unichain, because every state leads to $j$ in one step. Stationarity, $\varpi(\varepsilon)^{\top}\{(1-\varepsilon)P+\varepsilon\bm1\bm e_j^{\top}\}=\varpi(\varepsilon)^{\top}$, together with $\varpi(\varepsilon)^{\top}\bm1=1$, reads $\varpi(\varepsilon)^{\top}\{I-(1-\varepsilon)P\}=\varepsilon\bm e_j^{\top}$. The matrix in braces is invertible because the spectral radius of $(1-\varepsilon)P$ is $1-\varepsilon<1$. The path has direction $\Delta=\bm1\bm e_j^{\top}-P$, and $\varpi^{\top}\Delta=\bm e_j^{\top}-\varpi^{\top}$. By (a), $d\varpi_{\mathcal A}/d\varepsilon|_0=(\bm e_j-\varpi)^{\top}H\bm1_{\mathcal A}=g_{\mathcal A}(j)$, since $\varpi^{\top}H=\bm0^{\top}$. For every $y$, multiplying $H_{xy}=H_{yy}-\varpi_yM_{xy}$ by $\varpi_x$ and summing over $x$ gives $0=\varpi^{\top}H\bm e_y=H_{yy}-\varpi_y\sum_x\varpi_xM_{xy}$, so $H_{yy}=\varpi_y\E_\varpi[T_y]$. Hence $g_{\mathcal A}(j)=\sum_{y\in\mathcal A}H_{jy}=\sum_{y\in\mathcal A}(H_{yy}-\varpi_yM_{jy})=\sum_{y\in\mathcal A}\varpi_y(\E_\varpi[T_y]-M_{jy})$. Dividing by $\varpi(\mathcal A)>0$ gives the logarithmic derivative. If $\mathcal A=\{j\}$, then $M_{jj}=0$ and the expression reduces to $\E_\varpi[T_j]$. For $P=\bm1\varpi^{\top}$, $Z=(I-\bm1\varpi^{\top}+\bm1\varpi^{\top})^{-1}=I$ and $H=I-\bm1\varpi^{\top}$, so $g_{\mathcal A}(j)=1-\varpi(\mathcal A)$ for $j\in\mathcal A$.

\emph{(d)} Because $P$ is irreducible and $S\neq\mathcal X$, $q(S)>0$. Let $x\in S$. To reach $j\notin S$, the chain must first leave $S$. From each state of $S$ it stays in $S$ for one more step with probability at least $1-q(S)$, so by induction $\Pr_x(T_{\mathcal X\setminus S}>n)\ge\{1-q(S)\}^n$. Hence $M_{xj}\ge\E_xT_{\mathcal X\setminus S}=\sum_{n\ge0}\Pr_x(T_{\mathcal X\setminus S}>n)\ge1/q(S)$. By (b), $\kappa_j=\frac12\varpi_j\max_xM_{xj}\ge\varpi_j/\{2q(S)\}$, and $\E_\varpi[T_j]\ge\sum_{x\in S}\varpi_xM_{xj}\ge\varpi(S)/q(S)$.

\emph{(e)} The telescoping identity $\tilde P^t-P^t=\sum_{s=0}^{t-1}\tilde P^{s}\Delta P\,P^{t-1-s}$ and \eqref{eq:mk-eT} give
\[
  e_T(\tilde P,\nu)-e_T(P,\nu)=\frac1T\sum_{t=1}^{T-1}\sum_{s=0}^{t-1}(\nu^{\top}\tilde P^{s}\Delta P)\,(P^{t-1-s}\bm1_{\mathcal A}).
\]
Each row vector $\bm v^{\top}=\nu^{\top}\tilde P^s\Delta P$ sums to zero and has $\|\bm v\|_1\le\|\Delta P\|_\infty$, because $\nu^{\top}\tilde P^s$ is a probability vector. By \cref{lem:app-zerosum}(a) and (c), $|\bm v^{\top}P^r\bm1_{\mathcal A}|\le\frac12\|\Delta P\|_\infty\osc(P^r\bm1_{\mathcal A})\le\frac12\|\Delta P\|_\infty\operatorname{erg}(P^r)$. Summing over the pairs $(t,s)$, and noting that $r=t-1-s$ takes the value $r$ for $T-1-r$ pairs, gives the first bound. The second follows from $\operatorname{erg}(P^r)\le1$ and $\sum_{r=0}^{T-2}(T-1-r)=T(T-1)/2$. For the derivative, put $\Delta P=\varepsilon\Delta$, divide by $\varepsilon$ and let $\varepsilon\to0$, so that $\tilde P^s\to P^s$. \qed

\begin{remark}[Sharpness of \cref{thm:perturb}(b)]\label{rem:perturb-sharp}
Over formal directions, with $\Delta\bm1=\bm0$ and $\|\Delta\|_\infty\le1$, the supremum of $|d\varpi_{\mathcal A}/d\varepsilon|_{\varepsilon=0}$ equals $\kappa_{\mathcal A}$. It is attained by moving mass $\frac12$ in every row from a minimiser to a maximiser of $\bm g_{\mathcal A}$. Over directions that keep $P_{\mathrm{ref}}+\varepsilon\Delta$ stochastic for small $\varepsilon>0$, the largest first-order increase is $\frac12\sum_x\varpi_x\{\max_yg_{\mathcal A}(y)-\min_{y:P_{\mathrm{ref}}(x,y)>0}g_{\mathcal A}(y)\}$. This equals $\kappa_{\mathcal A}$ when every row of $P_{\mathrm{ref}}$ charges a minimiser of $\bm g_{\mathcal A}$, for instance when $P_{\mathrm{ref}}>0$, and is smaller otherwise.
The proof is the last paragraph of the proof of part (b) above.
\end{remark}

\subsection{Estimation from audit sessions}\label{app:markov-estimation}

\Cref{sec:markov-estimation} summarises this result. Within one arm of an audit, fix the policy and the persona and observe $n$ independent sessions $(X_{s1},\dots,X_{sT})$, $s=1,\dots,n$. Let $n_{0x}$ be the number of sessions that start in $x$, let $n_{xy}=\sum_s\sum_{t<T}\ind{X_{st}=x,X_{s,t+1}=y}$ count transitions, and let $n_{x+}=\sum_yn_{xy}$. Let $\bar n_x=\sum_{t=1}^{T-1}(\nu^{\top}P^{t-1})_x$ be the expected number of departures from $x$ per session, and $\Sigma_x=\operatorname{diag}(P_{x\cdot})-P_{x\cdot}^{\top}P_{x\cdot}$, where $P_{x\cdot}$ is the $x$th row of $P$ as a row vector. For a function $\Psi(\nu,P)$, let $\bm c=\partial\Psi/\partial\nu$ and let $G_{x\cdot}=\partial\Psi/\partial P_{x\cdot}$, a gradient determined up to a constant in each row.

\begin{proposition}[Estimation]\label{prop:markov-mle}
Let \cref{asm:markov} hold within the arm, and let $\bar n_x>0$ for every $x$.
\begin{enumerate}[(a)]
  \item The maximum-likelihood estimators are $\hat\nu_x=n_{0x}/n$ and $\hat P_{xy}=n_{xy}/n_{x+}$.
  \item As $n\to\infty$ with $T$ fixed, $\sqrt n(\hat\nu-\nu)$ and the rows $\sqrt n(\hat P_{x\cdot}-P_{x\cdot})$ are asymptotically independent and normal, with covariance matrices $\operatorname{diag}(\nu)-\nu\nu^{\top}$ and $\Sigma_x/\bar n_x$.
  \item If $\Psi$ is differentiable at $(\nu,P)$, then $\sqrt n\{\Psi(\hat\nu,\hat P)-\Psi(\nu,P)\}$ converges in distribution to $\mathrm N(0,\sigma^2_\Psi)$, with
  \[
    \sigma^2_\Psi=\Var_{x\sim\nu}(c_x)+\sum_x\frac{1}{\bar n_x}\Var_{y\sim P(x,\cdot)}(G_{xy}).
  \]
  \item For long-run exposure, $\Psi=\varpi(\mathcal A)$, one may take $\bm c=\bm0$ and $G_{xy}=\varpi_x\,g_{\mathcal A}(y)$, so that $\sigma^2_\Psi=\sum_x(\varpi_x^2/\bar n_x)\Var_{y\sim P(x,\cdot)}\{g_{\mathcal A}(y)\}$. For session exposure, $\Psi=e_T$, one may take $c_x=T^{-1}\sum_{t=0}^{T-1}(P^{t}\bm1_{\mathcal A})_x$ and $G_{xy}=T^{-1}\sum_{t=1}^{T-1}\sum_{s=0}^{t-1}(\nu^{\top}P^{s})_x(P^{t-1-s}\bm1_{\mathcal A})_y$.
\end{enumerate}
\end{proposition}

Parts (a) and (b) are the classical theory of inference for Markov chains. \Citet{anderson1957} treat exactly this design, many independent sequences of fixed length, and \citet{billingsley1961} surveys the case of one long chain. The likelihood factorises into one multinomial kernel per row, and each row behaves asymptotically like a multinomial sample of size $n\bar n_x$. Part (c) is the delta method \citep[Ch.~3]{vandervaart1998}, and part (d) supplies the gradients, the first from \cref{thm:perturb}(a). The two arms of an audit are independent, so the log amplification ratio $\log\{\hat\varpi_{\mathrm{alg}}(\mathcal A)/\hat\varpi_{\mathrm{ref}}(\mathcal A)\}$ has asymptotic variance $\sigma^2_{\mathrm{alg}}/\{n_{\mathrm{alg}}\varpi_{\mathrm{alg}}(\mathcal A)^2\}+\sigma^2_{\mathrm{ref}}/\{n_{\mathrm{ref}}\varpi_{\mathrm{ref}}(\mathcal A)^2\}$, with the $\sigma^2$ of part (d) for each arm.

\subsection{A sticky reference}\label{app:markov-nudge}

\Cref{tab:markov-nudge} isolates stickiness. Besides the reference $P_{\mathrm{ref}}$ of \cref{sec:markov-example}, consider the sticky reference
\[
  P_{\mathrm{stk}}=\begin{pmatrix}.994&.005&.001\\.05&.89&.06\\.05&.20&.75\end{pmatrix},
\]
with rows and columns ordered $(\mathsf N,\mathsf B,\mathsf A)$, in which neutral content rarely leads elsewhere and the borderline region is hard to leave. The two references have nearly the same long-run exposure, about $2.3\%$, and the same memoryless benchmark. Redirecting one recommendation in a hundred to $\mathsf A$ raises long-run exposure to $3.7\%$ under $P_{\mathrm{ref}}$, close to the memoryless $3.3\%$. Under the sticky reference it raises it to $8.0\%$, an increase of $5.6$ percentage points from a nudge of one percentage point. This is \cref{thm:perturb}(c) at work: the mean time to reach $\mathsf A$ from equilibrium is $262$ items under $P_{\mathrm{stk}}$ against $62$ under $P_{\mathrm{ref}}$. By \cref{thm:perturb}(d), with $S=\{\mathsf N\}$ and $q(S)=0.006$, it is at least $149$. The first-order prediction overshoots slightly, $8.5\%$, because a nudge of this size is no longer small for so sensitive a chain. Finite sessions show the same ordering. Under the sticky reference, $20$-item sessions see five times as much antisemitic content after the nudge as before.

\begin{table}[htb]
\centering
\caption{A nudge of size $\varepsilon=0.01$ toward $\mathsf A$ under two references with nearly the same long-run exposure (\cref{thm:perturb}(c)). \textbf{Illustrative, invented matrices; not data.} Sessions start in $\mathsf N$; $\mathcal A=\{\mathsf A\}$. All values are exact computations from the matrices in the text.}\label{tab:markov-nudge}
\small
\begin{tabular}{@{}l rr@{}}
\toprule
 & $P_{\mathrm{ref}}$ & $P_{\mathrm{stk}}$ \\
\midrule
Kemeny's constant $K$ & 3.13 & 21.1 \\
Sensitivity coefficient $\kappa_{\mathcal A}$ & 0.742 & 3.31 \\
Mean items to $\mathsf A$ from equilibrium, $\E_\varpi[T_{\mathsf A}]$ & 61.9 & 261.8 \\
Memoryless benchmark $(1-\varpi_{\mathsf A})/\varpi_{\mathsf A}$ & 42.1 & 41.3 \\
Long-run exposure $\varpi(\mathcal A)$: before nudge & 0.0232 & 0.0236 \\
\quad after nudge, exact & 0.0375 & 0.0796 \\
\quad after nudge, first order $\varpi_{\mathsf A}(1+\varepsilon\,\E_\varpi[T_{\mathsf A}])$ & 0.0376 & 0.0855 \\
\quad after nudge, memoryless $\varpi_{\mathsf A}+\varepsilon(1-\varpi_{\mathsf A})$ & 0.0330 & 0.0334 \\
Session exposure $e_{20}$: before nudge & 0.0207 & 0.0086 \\
\quad after nudge & 0.0339 & 0.0439 \\
Session exposure $e_{100}$: before nudge & 0.0227 & 0.0192 \\
\quad after nudge & 0.0368 & 0.0702 \\
\bottomrule
\end{tabular}
\end{table}

\subsection{Proof of \texorpdfstring{\cref{prop:markov-mle}}{the estimation proposition}}\label{app:pf-mle}

\emph{(a)} The likelihood of $n$ sessions is $\prod_s\nu_{X_{s1}}\prod_{t<T}P_{X_{st}X_{s,t+1}}=\prod_x\nu_x^{n_{0x}}\prod_x\prod_yP_{xy}^{n_{xy}}$. It is a product of multinomial kernels, one for $\nu$ and one for each row of $P$, and each is maximised by the empirical frequencies.

\emph{(b)} For session $s$, let $\bm Y^{\nu}_s=\bm e_{X_{s1}}-\nu$ and, for each $x$, $\bm Y_{s,x}=\sum_{t=1}^{T-1}\ind{X_{st}=x}(\bm e_{X_{s,t+1}}-P_{x\cdot}^{\top})$. These vectors are independent and identically distributed across sessions. Each summand of $\bm Y_{s,x}$ has conditional mean zero given $X_{s1},\dots,X_{st}$, so all the vectors have mean zero, and cross-products of summands at different times have mean zero. At equal times, $\ind{X_{st}=x}\ind{X_{st}=x'}=0$ for $x\neq x'$, and for $x=x'$ the conditional covariance of $\bm e_{X_{s,t+1}}$ given $X_{st}=x$ is $\Sigma_x$. Hence $\Cov(\bm Y_{s,x},\bm Y_{s,x'})=\ind{x=x'}\bar n_x\Sigma_x$. Conditioning on $X_{s1},\dots,X_{st}$ also shows that $\Cov(\bm Y^{\nu}_s,\bm Y_{s,x})=0$. The multivariate central limit theorem gives joint asymptotic normality of $n^{-1/2}\sum_s(\bm Y^\nu_s,(\bm Y_{s,x})_x)$, with block-diagonal covariance $\operatorname{diag}\{\operatorname{diag}(\nu)-\nu\nu^{\top},(\bar n_x\Sigma_x)_x\}$. Finally, $\sqrt n(\hat\nu-\nu)=n^{-1/2}\sum_s\bm Y^\nu_s$, and $\sqrt n(\hat P_{x\cdot}-P_{x\cdot})^{\top}=(n/n_{x+})\,n^{-1/2}\sum_s\bm Y_{s,x}$, where $n_{x+}/n\to\bar n_x>0$ almost surely by the strong law of large numbers. Slutsky's lemma completes the proof.

\emph{(c)} By the delta method \citep[Thm.~3.1]{vandervaart1998}, $\Psi(\hat\nu,\hat P)-\Psi(\nu,P)=\bm c^{\top}(\hat\nu-\nu)+\sum_xG_{x\cdot}(\hat P_{x\cdot}-P_{x\cdot})^{\top}+o_p(n^{-1/2})$. Its asymptotic variance is $\bm c^{\top}\{\operatorname{diag}(\nu)-\nu\nu^{\top}\}\bm c+\sum_x\bar n_x^{-1}G_{x\cdot}\Sigma_xG_{x\cdot}^{\top}$. Each quadratic form is a variance: $\bm c^{\top}\{\operatorname{diag}(\nu)-\nu\nu^{\top}\}\bm c=\Var_{x\sim\nu}(c_x)$ and $G_{x\cdot}\Sigma_xG_{x\cdot}^{\top}=\Var_{y\sim P(x,\cdot)}(G_{xy})$. Both are unchanged when a constant is added to $\bm c$ or to a row of $G$, which is why the gradients need only be determined up to such constants: estimation errors in $\hat\nu$ and in each row of $\hat P$ sum to zero.

\emph{(d)} $\varpi(\mathcal A)$ does not depend on $\nu$, so $\bm c=\bm0$. By \cref{thm:perturb}(a), the derivative of $\varpi(\mathcal A)$ along a direction $\Delta$ with $\Delta\bm1=\bm0$ is $\sum_x\varpi_x\sum_y\Delta_{xy}g_{\mathcal A}(y)=\sum_{x,y}G_{xy}\Delta_{xy}$ with $G_{xy}=\varpi_xg_{\mathcal A}(y)$. For $e_T$, the derivative along a change $\delta\nu$ of the initial distribution is $T^{-1}\sum_{t=0}^{T-1}\delta\nu^{\top}P^{t}\bm1_{\mathcal A}$. The derivative along $\Delta$ is the expression in \cref{thm:perturb}(e), $T^{-1}\sum_{t=1}^{T-1}\sum_{s=0}^{t-1}\nu^{\top}P^{s}\Delta P^{t-1-s}\bm1_{\mathcal A}=\sum_{x,y}G_{xy}\Delta_{xy}$, with $G$ as stated. \qed
